\documentclass[11pt,a4paper]{article}
\usepackage[T1]{fontenc}
\usepackage{isabelle,isabellesym}
\usepackage{pdfsetup}
\hypersetup{hidelinks}

\urlstyle{rm}
\isabellestyle{it}

\begin{document}

\title{Halting Absoluteness and Effective Sentence Invariance in Isabelle/ZF}
\author{Tang Ziyi}
\maketitle

\begin{abstract}
We construct a first-order set-theoretic formula whose satisfaction in every
transitive ZFC set model agrees with finite execution of a given set-coded
Turing machine. We then close the formula by naming its machine and input.
Let $I$ be the natural-number codes of closed formulas whose truth value
agrees across all such models. Assuming a countable transitive model of
ZFC exists, we construct primitive-recursive many-one reductions from both
self-input halting and its complement to $I$. The reductions use an
injective numerical code for formulas and the AFP formalisation of the
independence of the continuum hypothesis.
\end{abstract}

\tableofcontents

\section{Halting formula}
For every machine $P$, input $x$, and transitive set model $M$ of ZFC, the
explicit formula $\mathit{halt\_fm}$ satisfies
\[
  (M,[P,x])\models\mathit{halt\_fm}
  \quad\Longleftrightarrow\quad
  \mathrm{halts}(P,x).
\]
Theorem \texttt{transitive\_model\_halt\_fm\_iff} proves this statement.
Finite run witnesses and their closure in transitive models are established
before the formula is introduced. The closed sentence
$\mathit{halt\_sentence}(P,x)$ names both parameters internally, and
\texttt{sats\_halt\_sentence\_iff} proves its corresponding satisfaction
equation. These results do not use CH or an assumption that a countable
transitive model exists.

\section{Effective reduction}
Let $K$ be the self-input halting set of the companion computation entry,
and let $I$ contain exactly the numerical codes of closed formulas whose
truth value agrees across all transitive set models of ZFC. Given a
countable transitive set model of ZFC, this entry constructs
primitive-recursive $f_+$ and $f_-$ such that, for every natural number $e$,
\[
  e\in K\ \Longleftrightarrow\ f_+(e)\in I,
  \qquad
  e\notin K\ \Longleftrightarrow\ f_-(e)\in I.
\]
Theorems \texttt{self\_halting\_pr\_many\_one\_invariance} and
\texttt{nonhalting\_pr\_many\_one\_invariance} are the formal reduction
statements. CH supplies a fixed sentence on which two transitive models
disagree, using the AFP formalisation of its independence~\cite{afp-ch}.
The countable-model assumption is needed for these disagreement witnesses.

The target set contains only codes of well-formed closed formulas. A
constructor-tagged pairing code is proved injective on formulas. The two
reducing functions are proved members of Isabelle/ZF's class
\texttt{prim\_rec}, with application equations identifying their values
with the codes of the intended sentences. Their finite-list component uses
a bounded reverse traversal of the machine code.

The machine semantics and self-input halting theorem come from the companion
entry~\cite{afp-core}. The entry proves primitive-recursive many-one
reductions. It does not identify every primitive-recursive function with a
machine in that specific Turing model.

\input{session}

\section*{Development provenance}
Tang Ziyi formulated the research question and wrote the initial
formalisation. Subsequent definitions, proofs, and repository engineering
include substantial assistance from OpenAI Codex.

\begin{thebibliography}{9}

\bibitem{afp-ch}
Emmanuel Gunther, Miguel Pagano, Pedro S\'anchez Terraf, and Mat\'ias Steinberg.
\emph{The Independence of the Continuum Hypothesis in Isabelle/ZF}.
Archive of Formal Proofs, 2022.
\url{https://isa-afp.org/entries/Independence_CH.html}.

\bibitem{afp-core}
Tang Ziyi.
\emph{Set-Coded Computation in Isabelle/ZF}.
Companion formal development, source repository.

\end{thebibliography}

\end{document}
